% !TEX root = ../jalase03.tex
%% ============================================================
%%  جلسهٔ ۳ — نتایج اصول کولموگروف و مدل‌های احتمال گسسته
%%  (احتمال شرطی، احتمال کل و بیز: sessions/session04.tex)
%% ============================================================
\chapter{نتایج اصول کولموگروف و مدل‌های احتمال گسسته}
\label{ch:s3}

\sessioncard{شنبه ۱۱ مهر}
{ادامهٔ فضای احتمال: نتایج اصول کولموگروف و مدل‌های گسسته}
{\lr{Pishro-Nik} بخش‌های \lr{1.3} و \lr{2.1} \,؛\, \lr{Ross} بخش‌های \lr{2.4} و \lr{2.5} \,؛\, \lr{Bertsekas} بخش \lr{1.2}}
{\begin{itemize}
  \item قواعد پایهٔ احتمال (متمم، یکنوایی، تفاضل) را فقط از سه اصل کولموگروف \emph{اثبات} کنید؛
  \item برای فضای نمونهٔ متناهی یا شمارا مدل احتمال بسازید و مسائل «هم‌شانس» را با ترکیبیات حل کنید؛
  \item اصل شمول و عدم شمول را اثبات کنید و در مسائل شمارشی و گراف تصادفی به کار ببرید؛
  \item کران اجتماع و نامساوی بونفرونی را اثبات و در تحلیل قابلیت اطمینان و «روش احتمالاتی» استفاده کنید؛
  \item پیوستگی احتمال را برای دنباله‌های یکنوا بیان و اثبات کنید.
\end{itemize}}

%% ------------------------------------------------------------
\section{مرور: فضای احتمال و سه اصل کولموگروف}
%% ------------------------------------------------------------

در جلسهٔ قبل دیدیم که یک مدل احتمال یعنی سه‌تایی $(\Omega,\F,\P)$: فضای نمونهٔ $\Omega$، یک $\sigma$-میدان $\F$ از پیشامدها، و قانون احتمال $\P$ که به هر پیشامد وزنی می‌دهد و سه اصل کولموگروف را برآورده می‌کند:

\begin{keyformula}
\[ \P(A)\ge0,\qquad \P(\Omega)=1,\qquad A_i\cap A_j=\varnothing\ (i\ne j)\ \Rightarrow\ \P\Bigl(\bigcup_i A_i\Bigr)=\sum_i\P(A_i) \]
\end{keyformula}

\begin{intuition}
چرا این اصول معقول‌اند؟ فراوانی نسبی را در نظر بگیرید. اگر آزمایش $N$ بار تکرار شود:
$N_A/N\ge0$ (اصل ۱)، $\;N_\Omega/N=N/N=1$ چون $\Omega$ همیشه رخ می‌دهد (اصل ۲)، و اگر $A$ و $B$ هم‌زمان رخ ندهند، $N_{A\cup B}=N_A+N_B$ (اصل ۳ برای دو پیشامد). یعنی اصول کولموگروف دقیقاً ویژگی‌هایی را می‌خواهند که فراوانی نسبی به طور طبیعی دارد.
\end{intuition}

\begin{warning}
اصل ۳ فقط برای پیشامدهای \emph{مجزا} است. در حالت کلی $\P(A\cup B)\ne\P(A)+\P(B)$. مثلاً در پرتاب یک تاس، اگر $A=\{2,4,6\}$ و $B=\{4,5,6\}$ آنگاه $\P(A)+\P(B)=1$ ولی $\P(A\cup B)=\P(\{2,4,5,6\})=\tfrac46$.
\end{warning}

در این جلسه دو کار می‌کنیم: (۱) نشان می‌دهیم \emph{همهٔ} قواعد آشنای احتمال فقط از همین سه اصل نتیجه می‌شوند؛ (۲) مدل‌های \emph{گسسته} را، که در آن‌ها $\F=2^\Omega$ است، به‌طور کامل می‌سازیم و با مثال‌های متنوع به کار می‌بریم.

%% ------------------------------------------------------------
\section{قواعد پایه: نتایج مستقیم اصول}
%% ------------------------------------------------------------

همهٔ قواعد آشنای احتمال از همین سه اصل نتیجه می‌شوند. در ادامه $(\Omega,\F,\P)$ یک فضای احتمال و همهٔ مجموعه‌ها عضو $\F$اند.

\begin{theorem}[label=thm:basic-props]{خواص پایه}
\begin{enumerate}
  \item $\P(\varnothing)=0$.
  \item \textbf{جمع‌پذیری متناهی:} اگر $A_1,\dots,A_n$ دوبه‌دو مجزا باشند، $\P\bigl(\bigcup_{i=1}^nA_i\bigr)=\sum_{i=1}^n\P(A_i)$.
  \item \textbf{قاعدهٔ متمم:} $\P(A^c)=1-\P(A)$.
  \item $0\le\P(A)\le1$.
  \item \textbf{قاعدهٔ تفاضل:} $\P(A-B)=\P(A)-\P(A\cap B)$.
  \item \textbf{یکنوایی:} اگر $A\subseteq B$ آنگاه $\P(A)\le\P(B)$ و $\P(B-A)=\P(B)-\P(A)$.
\end{enumerate}
\end{theorem}

\begin{proof}
\begin{enumerate}
  \item دنبالهٔ $A_1=A_2=\cdots=\varnothing$ دوبه‌دو مجزاست و اجتماعش $\varnothing$ است. طبق اصل ۳: $\P(\varnothing)=\sum_{i=1}^\infty\P(\varnothing)$. اگر $\P(\varnothing)=c>0$ بود، طرف راست $+\infty$ می‌شد که ممکن نیست ($\P$ مقدار حقیقی دارد). پس با توجه به اصل ۱، $\P(\varnothing)=0$.
  \item قرار دهید $A_{n+1}=A_{n+2}=\cdots=\varnothing$ و اصل ۳ و بند (۱) را به کار ببرید.
  \item $A$ و $A^c$ مجزا هستند و $A\cup A^c=\Omega$. طبق بند (۲) و اصل ۲: $\P(A)+\P(A^c)=\P(\Omega)=1$.
  \item $\P(A)\ge0$ اصل ۱ است و چون $\P(A^c)\ge0$، از بند (۳) داریم $\P(A)=1-\P(A^c)\le1$.
  \item طبق قضیهٔ «شکستن مجموعه با افراز» (جلسهٔ ۱)، $A=(A\cap B)\cup(A-B)$ و این دو مجزا هستند؛ پس $\P(A)=\P(A\cap B)+\P(A-B)$.
  \item اگر $A\subseteq B$ آنگاه $A\cap B=A$ و طبق بند (۵): $\P(B-A)=\P(B)-\P(A)$. چون $\P(B-A)\ge0$، نتیجه می‌شود $\P(A)\le\P(B)$.
\end{enumerate}
\end{proof}

\begin{example}{یکنوایی در عمل}
یک سکهٔ سالم را $n\ge10$ بار پرتاب می‌کنیم. فرض کنید $A$ پیشامد «دست‌کم ۱۰ شیر بیاید» و $B$ پیشامد «۱۰ پرتاب اول همه شیر باشند» باشد. چون $B\subseteq A$، بدون هیچ محاسبهٔ پیچیده‌ای می‌دانیم $\P(A)\ge\P(B)$. به همین ترتیب اگر بدانیم $\P(A)=0.01$، مطمئنیم $\P(B)\le0.01$.
\end{example}

%% ------------------------------------------------------------
\section{مدل‌های احتمال گسسته}
\label{sec:discrete}
%% ------------------------------------------------------------

\begin{definition}{فضای احتمال گسسته}
اگر $\Omega$ متناهی یا شمارا باشد و $\F=2^\Omega$، فضای احتمال را \term{گسسته}{discrete} می‌نامیم. تابع
\[ p(\omega)\defeq\P(\{\omega\}),\qquad \omega\in\Omega \]
را \emph{تابع جرم} (وزنِ) برآمدها می‌گوییم.
\end{definition}

\begin{theorem}[label=thm:discrete]{ساختن احتمال گسسته از وزن‌ها}
\begin{enumerate}
  \item در هر فضای احتمال گسسته، برای هر پیشامد $A$:
  \[ \P(A)=\sum_{\omega\in A}p(\omega),\qquad\text{و به‌ویژه}\qquad\sum_{\omega\in\Omega}p(\omega)=1. \]
  \item برعکس، اگر $\Omega$ شمارا باشد و $p:\Omega\to[0,\infty)$ تابعی با $\sum_{\omega\in\Omega}p(\omega)=1$ باشد، آنگاه $\P(A)\defeq\sum_{\omega\in A}p(\omega)$ یک قانون احتمال روی $2^\Omega$ است.
\end{enumerate}
\end{theorem}

\begin{proof}
(۱) $A$ شماراست، پس $A=\bigcup_{\omega\in A}\{\omega\}$ اجتماعی شمارا از پیشامدهای سادهٔ دوبه‌دو مجزاست. طبق اصل ۳، $\P(A)=\sum_{\omega\in A}\P(\{\omega\})$. با $A=\Omega$ و اصل ۲ رابطهٔ دوم به دست می‌آید.\\
(۲) اصل ۱ از $p\ge0$ و اصل ۲ از فرض $\sum p=1$ نتیجه می‌شود. برای اصل ۳: اگر $A_1,A_2,\dots$ مجزا باشند، $\sum_{\omega\in\bigcup A_i}p(\omega)=\sum_i\sum_{\omega\in A_i}p(\omega)$؛ این جابه‌جایی و دسته‌بندی جمله‌ها مجاز است چون سری جمله‌های نامنفی دارد (ترتیب جمع در سری‌های نامنفی اهمیتی ندارد).
\end{proof}

\begin{intuition}
در حالت گسسته کل «جرم احتمال» روی نقاط منفرد نشسته است: کافی است وزن هر برآمد را بدانیم؛ احتمال هر پیشامد جمع وزن‌های نقاط آن است. این همان کاری است که در جلسهٔ ۵ با \emph{تابع جرم احتمال} (\lr{PMF}) متغیرهای تصادفی گسسته انجام می‌دهیم.
\end{intuition}

\subsection{مدل کلاسیک: برآمدهای هم‌شانس}

\begin{corollary}[label=cor:classical]{مدل هم‌شانس}
اگر $\Omega$ متناهی باشد و همهٔ برآمدها هم‌شانس باشند، یعنی $p(\omega)=c$ برای هر $\omega$، آنگاه $c=1/|\Omega|$ و
\[ \P(A)=\frac{|A|}{|\Omega|}. \]
\end{corollary}

\begin{proof}
طبق قضیهٔ \ref{thm:discrete}: $1=\sum_{\omega\in\Omega}c=c\,|\Omega|$، پس $c=1/|\Omega|$ و $\P(A)=\sum_{\omega\in A}c=|A|/|\Omega|$.
\end{proof}

پس تعریف کلاسیکِ لاپلاس (جلسهٔ ۱) دیگر یک «تعریف» نیست، بلکه \emph{حالت خاصی از اصول کولموگروف} است.

\begin{warning}
روی یک فضای نمونهٔ \emph{نامتناهیِ شمارا} مدل هم‌شانس وجود ندارد. مثلاً «یک عدد طبیعی را کاملاً تصادفی انتخاب کن» معنای دقیقی ندارد: اگر $p(n)=c$ برای همهٔ $n\in\N$، آنگاه $\sum_np(n)$ برای $c=0$ برابر $0$ و برای $c>0$ برابر $\infty$ است و هیچ‌کدام ۱ نیست.
\end{warning}

\begin{history}{اشتباه لایبنیتس}
گوتفرید لایبنیتس، یکی از بزرگ‌ترین ریاضی‌دانان تاریخ، معتقد بود با دو تاس، گرفتن مجموع ۱۱ و مجموع ۱۲ هم‌شانس است، چون هر کدام «فقط یک راه» دارند: $\{5,6\}$ و $\{6,6\}$. اشتباه او این بود که برآمدهای \emph{بدون ترتیب} را هم‌شانس فرض کرد. با تاس‌های قابل تشخیص (مثلاً یکی سفید و یکی قرمز)، مجموع ۱۱ دو برآمد $(5,6)$ و $(6,5)$ دارد ولی مجموع ۱۲ فقط یکی: $\P(11)=2/36$ ولی $\P(12)=1/36$.
\emph{درس مهم:} در مدل هم‌شانس، فضای نمونه را طوری بسازید که برآمدها \emph{واقعاً} هم‌شانس باشند؛ معمولاً یعنی اشیا را متمایز و ترتیب را مهم در نظر بگیرید.
\end{history}

\begin{example}[label=ex:n-coins]{$n$ پرتاب سکه}
یک سکهٔ سالم $n$ بار پرتاب می‌شود. احتمال اینکه دقیقاً $k$ شیر بیاید چقدر است؟
\begin{solution}
$\Omega=\{H,T\}^n$ با $2^n$ برآمد هم‌شانس (هر دنباله احتمال $2^{-n}$ دارد). تعداد دنباله‌هایی که دقیقاً $k$ شیر دارند برابر تعداد روش‌های انتخاب \emph{جایگاه} شیرهاست، یعنی $\binom nk$. پس
\[ \P(\fa{دقیقاً $k$ شیر})=\binom nk2^{-n},\qquad k=0,1,\dots,n. \]
بررسی سازگاری: طبق اتحاد $\sum_k\binom nk=2^n$ از جلسهٔ ۱، $\sum_{k=0}^n\binom nk2^{-n}=2^n\cdot2^{-n}=1$؛ یا با قضیهٔ دوجمله‌ای: $\sum_k\binom nk(\tfrac12)^k(\tfrac12)^{n-k}=(\tfrac12+\tfrac12)^n=1$.
همچنین $\P(\fa{دست‌کم $k$ شیر})=\sum_{i=k}^n\binom ni2^{-n}$. مثلاً برای $n=10$، احتمال دست‌کم ۸ شیر برابر است با $(45+10+1)/1024=56/1024\approx0.055$.
(این همان \emph{توزیع دوجمله‌ای} است که در جلسهٔ ۵ به تفصیل بررسی می‌کنیم.)
\end{solution}
\end{example}

\begin{example}{پرتاب تا اولین شیر}
سکهٔ سالمی را آن‌قدر پرتاب می‌کنیم تا اولین شیر ظاهر شود. (الف) یک مدل احتمال بسازید. (ب) احتمال اینکه تعداد پرتاب‌ها زوج باشد چقدر است؟
\begin{solution}
(الف) $\Omega=\{H,\ TH,\ TTH,\ \dots\}$ که شماراست؛ برآمدِ $\omega_k=\underbrace{T\cdots T}_{k-1}H$ یعنی «اولین شیر در پرتاب $k$-ام». طبیعی است وزن $p(\omega_k)=2^{-k}$ را بدهیم (احتمال یک دنبالهٔ مشخص به طول $k$). این یک مدل معتبر است، چون $\sum_{k=1}^\infty2^{-k}=1$ (سری هندسی).
(ب) $\P(\text{زوج})=\sum_{j=1}^\infty p(\omega_{2j})=\sum_{j=1}^\infty4^{-j}=\dfrac{1/4}{1-1/4}=\dfrac13.$
(مدل این مثال همان \emph{توزیع هندسی} جلسهٔ ۶ است.)
\end{solution}
\end{example}

\begin{example}[label=ex:birthday]{مسئلهٔ روز تولد}
در کلاسی با $n$ دانشجو، احتمال اینکه دست‌کم دو نفر روز تولد یکسان داشته باشند چقدر است؟ (سال را ۳۶۵ روزه و همهٔ روزها را هم‌شانس فرض کنید.)
\begin{solution}
فضای نمونه همهٔ $n$تایی‌های مرتب از روزهاست: $|\Omega|=365^n$. محاسبهٔ مستقیمِ «دست‌کم یک تکرار» دشوار است، پس از متمم استفاده می‌کنیم. تعداد $n$تایی‌هایی که همه متمایزند $365\cdot364\cdots(365-n+1)$ است؛ بنابراین
\[ \P(\text{دست‌کم یک تکرار})=1-\frac{365\cdot364\cdots(365-n+1)}{365^n}=1-\prod_{i=0}^{n-1}\Bigl(1-\frac{i}{365}\Bigr). \]
\begin{figblock}
\begin{tikzpicture}
\begin{axis}[
  width=11cm, height=5.4cm, axis lines=left,
  xmin=0, xmax=72, ymin=0, ymax=1.05,
  xtick={0,10,23,30,40,50,60,70}, ytick={0,0.25,0.5,0.75,1},
  tick label style={font=\footnotesize}, grid=major, grid style={graytext!20},
  xlabel={$n$}, every axis x label/.style={at={(ticklabel* cs:1.02)}, anchor=west}]
  \addplot[thick, accent, smooth] coordinates {(1,0.0000) (4,0.0164) (7,0.0562) (10,0.1169) (13,0.1944) (16,0.2836) (19,0.3791) (22,0.4757) (25,0.5687) (28,0.6545) (31,0.7305) (34,0.7953) (37,0.8487) (40,0.8912) (43,0.9239) (46,0.9483) (49,0.9658) (52,0.9780) (55,0.9863) (58,0.9917) (61,0.9951) (64,0.9972) (67,0.9984) (70,0.9992)};
  \draw[dashed, primary] (axis cs:23,0) -- (axis cs:23,0.5073) -- (axis cs:0,0.5073);
  \fill[primary] (axis cs:23,0.5073) circle (2pt);
  \node[font=\footnotesize, primary, anchor=west] at (axis cs:25,0.44) {$n=23:\ 0.507$};
\end{axis}
\end{tikzpicture}

\figcaption[fig:birthday]{احتمال وجود دست‌کم یک روز تولد مشترک بر حسب تعداد افراد.}
\end{figblock}
\begin{center}
\small
\begin{tabular}{c|ccccccc}
$n$ & $10$ & $20$ & $23$ & $30$ & $40$ & $50$ & $70$\\\hline
احتمال & $0.117$ & $0.411$ & $0.507$ & $0.706$ & $0.891$ & $0.970$ & $0.999$
\end{tabular}
\end{center}
فقط با ۲۳ نفر احتمال بیش از ۵۰٪ است! شهود ما خطا می‌کند چون به «تکرار با \emph{روز تولد من}» فکر می‌کنیم، در حالی که $\binom{23}{2}=253$ \emph{زوج} مختلف وجود دارد که هر کدام می‌توانند تکرار ایجاد کنند. با تقریب $1-x\approx e^{-x}$:
\[ \P(\text{دست‌کم یک تکرار})\approx1-e^{-n(n-1)/730}. \]
\end{solution}
\end{example}

\begin{example}{معمای شوالیه دومِره}
کدام محتمل‌تر است: (الف) در ۴ پرتاب یک تاس دست‌کم یک ۶ بیاید؛ (ب) در ۲۴ پرتاب یک جفت تاس دست‌کم یک جفت‌شش بیاید؟
\begin{solution}
(الف) از $6^4$ برآمد هم‌شانس، $5^4$ برآمد هیچ ۶ ندارند؛ پس احتمال برابر $1-(5/6)^4\approx0.518$ است.
(ب) هر پرتاب جفت تاس ۳۶ برآمد دارد که ۳۵تای آن جفت‌شش نیست؛ پس احتمال برابر $1-(35/36)^{24}\approx0.491$ است.
دومِره (به‌درستی!) از تجربه‌اش در قمار دیده بود که (الف) کمی بیشتر رخ می‌دهد، و این تناقضی با استدلال سادهٔ «نسبت‌ها یکسان است» ($4/6=24/36$) داشت. پاسکال با همین محاسبه معما را حل کرد.
\end{solution}
\end{example}

%% ------------------------------------------------------------
\section{اصل شمول و عدم شمول}
%% ------------------------------------------------------------


\begin{theorem}[label=thm:ie2]{شمول و عدم شمول برای دو پیشامد}
برای هر دو پیشامد $A$ و $B$:
\[ \P(A\cup B)=\P(A)+\P(B)-\P(A\cap B). \]
\end{theorem}

\begin{proof}
$A\cup B$ را به دو قطعهٔ مجزا بشکنید: $A\cup B=A\cup(B-A)$. پس
\[ \P(A\cup B)=\P(A)+\P(B-A)=\P(A)+\P(B)-\P(A\cap B), \]
که در تساوی دوم از قاعدهٔ تفاضل استفاده شد. (شهود: در $\P(A)+\P(B)$ ناحیهٔ مشترک دو بار شمرده شده است؛ یک بار کمش می‌کنیم.)
\end{proof}

\begin{theorem}[label=thm:ie-general]{اصل شمول و عدم شمول (حالت کلی)}
برای پیشامدهای $A_1,\dots,A_n$:
\[
\P\Bigl(\bigcup_{i=1}^n A_i\Bigr)=\sum_{i}\P(A_i)-\sum_{i<j}\P(A_i\cap A_j)+\sum_{i<j<k}\P(A_i\cap A_j\cap A_k)-\cdots+(-1)^{n+1}\P(A_1\cap\cdots\cap A_n).
\]
به بیان فشرده: $\displaystyle \P\Bigl(\bigcup_{i=1}^n A_i\Bigr)=\sum_{k=1}^{n}(-1)^{k+1}S_k$ که در آن $\displaystyle S_k=\sum_{i_1<\cdots<i_k}\P(A_{i_1}\cap\cdots\cap A_{i_k})$.
\end{theorem}

\begin{proof}
با استقرا روی $n$. حالت $n=1$ بدیهی و حالت $n=2$ قضیهٔ \ref{thm:ie2} است. فرض کنید حکم برای $n$ پیشامد درست باشد و قرار دهید $B=\bigcup_{i=1}^nA_i$. طبق حالت دوتایی:
\[ \P\Bigl(\bigcup_{i=1}^{n+1}A_i\Bigr)=\P(B)+\P(A_{n+1})-\P(B\cap A_{n+1}),\qquad B\cap A_{n+1}=\bigcup_{i=1}^n(A_i\cap A_{n+1}). \]
فرض استقرا را یک بار برای $\P(B)$ و یک بار برای اجتماعِ $n$ پیشامدِ $A_i\cap A_{n+1}$ به کار ببرید. جمله‌های $\P(B)$ همهٔ اشتراک‌هایی‌اند که $A_{n+1}$ را ندارند، و جمله‌های $-\P(B\cap A_{n+1})$ (با علامت‌های برعکس‌شده) همهٔ اشتراک‌های شامل $A_{n+1}$ به جز خودِ $\P(A_{n+1})$اند. با افزودن $\P(A_{n+1})$، هر اشتراکِ $k$تایی از $\{A_1,\dots,A_{n+1}\}$ دقیقاً یک بار با علامت $(-1)^{k+1}$ ظاهر می‌شود.
\end{proof}

مثال زیر کاربرد جذابی از این اصل است که جوابش، به شکلی غیرمنتظره، عدد $e$ را نشان می‌دهد.

\begin{example}[label=ex:matching]{مسئلهٔ تطابق (کلاه‌ها)}
$n$ نفر کلاه‌هایشان را به رختکن می‌دهند و کلاه‌ها به صورت کاملاً تصادفی (همهٔ $n!$ جایگشت هم‌شانس) به آن‌ها برگردانده می‌شود. احتمال اینکه دست‌کم یک نفر کلاه خودش را بگیرد چقدر است؟ وقتی $n\to\infty$ چه می‌شود؟
\begin{solution}
فرض کنید $A_i$ پیشامد «نفر $i$ کلاه خودش را بگیرد» باشد. برای هر $i_1<\cdots<i_k$، در پیشامد $A_{i_1}\cap\cdots\cap A_{i_k}$ این $k$ نفر کلاه خود را می‌گیرند و بقیه به هر ترتیبی: $(n-k)!$ جایگشت. پس
\[ \P(A_{i_1}\cap\cdots\cap A_{i_k})=\frac{(n-k)!}{n!},\qquad S_k=\binom nk\frac{(n-k)!}{n!}=\frac1{k!}. \]
طبق اصل شمول و عدم شمول (قضیهٔ \ref{thm:ie-general}):
\[ \P\Bigl(\bigcup_{i=1}^nA_i\Bigr)=\sum_{k=1}^n\frac{(-1)^{k+1}}{k!}=1-\frac1{2!}+\frac1{3!}-\cdots\ \xrightarrow{\ n\to\infty\ }\ 1-e^{-1}\approx0.632. \]
جالب اینکه این احتمال تقریباً به $n$ بستگی ندارد: برای $n=5$ برابر $0.633$ و برای $n=7$ برابر $0.6321$ است!
\end{solution}
\end{example}

%% ------------------------------------------------------------
\section{کران اجتماع و نامساوی بونفرونی}
%% ------------------------------------------------------------


\begin{theorem}[label=thm:union-bound]{کران اجتماع (نامساوی بول)}
برای هر دنبالهٔ (متناهی یا شمارا) از پیشامدها:
\[ \P\Bigl(\bigcup_{i}A_i\Bigr)\le\sum_{i}\P(A_i). \]
\end{theorem}

\begin{proof}
پیشامدها را «مجزاسازی» می‌کنیم: $B_1=A_1$ و برای $i\ge2$، $B_i=A_i-\bigcup_{j<i}A_j$ (یعنی بخشی از $A_i$ که در $A_j$های قبلی نبوده). $B_i$ها دوبه‌دو مجزا هستند، $B_i\subseteq A_i$، و $\bigcup_iB_i=\bigcup_iA_i$ (هر نقطهٔ اجتماع، در $B_i$ متناظر با \emph{اولین} $A_i$ای که شاملش است قرار می‌گیرد). پس طبق اصل ۳ و یکنوایی:
\[ \P\Bigl(\bigcup_iA_i\Bigr)=\P\Bigl(\bigcup_iB_i\Bigr)=\sum_i\P(B_i)\le\sum_i\P(A_i). \]
\end{proof}

\begin{corollary}[label=cor:bonferroni]{نامساوی بونفرونی}
\[ \P(A\cap B)\ge\P(A)+\P(B)-1,\qquad \P\Bigl(\bigcap_{i=1}^nA_i\Bigr)\ge1-\sum_{i=1}^n\P(A_i^c). \]
\end{corollary}

\begin{proof}
طبق دمورگان و کران اجتماع: $\P\bigl(\bigcap_iA_i\bigr)=1-\P\bigl(\bigcup_iA_i^c\bigr)\ge1-\sum_i\P(A_i^c)$. حالت اول، حالت خاص $n=2$ است.
\end{proof}

\begin{example}{قابلیت اطمینان یک سیستم}
یک سیستم از ۲۰ قطعه تشکیل شده و فقط وقتی کار می‌کند که \emph{همهٔ} قطعات سالم باشند. احتمال خرابی هر قطعه در طول یک سال حداکثر $0.001$ است. بدون دانستن هیچ اطلاعاتی دربارهٔ ارتباط خرابی‌ها با هم، دربارهٔ احتمال سالم ماندن سیستم چه می‌توان گفت؟
\begin{solution}
اگر $A_i$ پیشامد «قطعهٔ $i$ سالم بماند» باشد، طبق بونفرونی
$\P(\text{سیستم سالم})=\P\bigl(\bigcap_{i=1}^{20}A_i\bigr)\ge1-20\times0.001=0.98$.
نکتهٔ مهم این است که این کران \emph{همیشه} درست است، حتی اگر خرابی قطعات به هم وابسته باشند.
\end{solution}
\end{example}

%% ------------------------------------------------------------
\section{پیوستگی احتمال}
%% ------------------------------------------------------------


\begin{theorem}[label=thm:continuity]{پیوستگی احتمال}
\begin{enumerate}
  \item اگر $A_1\subseteq A_2\subseteq\cdots$ (دنبالهٔ \emph{صعودی}) آنگاه $\displaystyle\P\Bigl(\bigcup_{n=1}^\infty A_n\Bigr)=\lim_{n\to\infty}\P(A_n)$.
  \item اگر $A_1\supseteq A_2\supseteq\cdots$ (دنبالهٔ \emph{نزولی}) آنگاه $\displaystyle\P\Bigl(\bigcap_{n=1}^\infty A_n\Bigr)=\lim_{n\to\infty}\P(A_n)$.
\end{enumerate}
\end{theorem}

\begin{example}{سکه‌ای که هرگز شیر نمی‌آید}
سکهٔ سالمی را بی‌نهایت بار پرتاب می‌کنیم. احتمال اینکه \emph{هرگز} شیر نیاید چقدر است؟
\begin{solution}
فرض کنید $A_n$ پیشامد «$n$ پرتاب اول همه خط باشند» باشد. از $2^n$ دنبالهٔ هم‌شانس برای $n$ پرتاب اول فقط یکی همه‌خط است؛ پس $\P(A_n)=2^{-n}$. دنبالهٔ $A_n$ نزولی است و $\bigcap_nA_n$ همان پیشامد «هرگز شیر نیاید» است. طبق پیوستگی:
$\P(\text{هرگز شیر نیاید})=\lim_{n\to\infty}2^{-n}=0$.
\end{solution}
\end{example}

\begin{note}
پیشامدی با احتمال صفر لزوماً «ناممکن» نیست: دنبالهٔ «همه خط» یک برآمد کاملاً مجاز است، ولی احتمالش صفر است. به همین ترتیب پیشامدی با احتمال ۱ را «تقریباً حتمی» (\lr{almost sure}) می‌نامیم.
\end{note}

%% ------------------------------------------------------------
\section*{جمع‌بندی جلسهٔ سوم}
\addcontentsline{toc}{section}{خلاصه و تمرین‌ها}
%% ------------------------------------------------------------

\begin{summary}
\begin{itemize}
  \item همهٔ قواعد احتمال از \textbf{سه اصل کولموگروف} نتیجه می‌شوند: $\P(\varnothing)=0$، جمع‌پذیری متناهی، $\P(A^c)=1-\P(A)$، $\P(A)\le1$، $\P(A-B)=\P(A)-\P(A\cap B)$ و \textbf{یکنوایی}.
  \item \textbf{مدل گسسته:} اگر $\Omega$ شمارا باشد، $\F=2^\Omega$ و $\P(A)=\sum_{\omega\in A}p(\omega)$ با $\sum_\omega p(\omega)=1$. در مدل \textbf{هم‌شانس}: $\P(A)=|A|/|\Omega|$؛ روی فضای شمارای نامتناهی مدل هم‌شانس وجود ندارد.
  \item در مسائل هم‌شانس فضای نمونه را با اشیای \emph{متمایز} بسازید (درس اشتباه لایبنیتس) و برای «دست‌کم یکی» از \emph{متمم} استفاده کنید (مسئلهٔ روز تولد، دومِره).
  \item \textbf{شمول و عدم شمول:} $\P(A\cup B)=\P(A)+\P(B)-\P(A\cap B)$ و تعمیم آن با علامت‌های یک‌درمیان؛ کاربرد: مسئلهٔ کلاه‌ها ($\to1-1/e$).
  \item \textbf{کران اجتماع:} $\P(\bigcup A_i)\le\sum\P(A_i)$ و \textbf{بونفرونی:} $\P(\bigcap A_i)\ge1-\sum\P(A_i^c)$؛ بدون نیاز به دانستن وابستگی پیشامدها. کاربرد: قابلیت اطمینان سیستم‌ها.
  \item \textbf{پیوستگی:} برای دنباله‌های صعودی یا نزولی، احتمالِ حد برابر حدِ احتمال‌هاست. پیشامد با احتمال صفر لزوماً ناممکن نیست.
\end{itemize}
\end{summary}



\subsection*{تمرین‌های پیشنهادی}

\begin{exercise}
با استفاده از اصول موضوعه ثابت کنید: (الف) $\P(A\cup B)\le\P(A)+\P(B)$؛ \ (ب) $\P(A\,\triangle\,B)=\P(A)+\P(B)-2\P(A\cap B)$؛ \ (ج) اگر $\P(A)=\P(B)=1$ آنگاه $\P(A\cap B)=1$.
\end{exercise}

\begin{exercise}
اگر $\P(A)=0.6$، $\P(B)=0.5$ و $\P(A\cup B)=0.8$، مقادیر $\P(A\cap B)$، $\P(A\cap B^c)$ و $\P(A^c\cap B^c)$ را بیابید.
\answer{$0.3$، \ $0.3$، \ $0.2$}
\end{exercise}

\begin{exercise}
فرض کنید $\P(A)=0.9$ و $\P(B)=0.8$. کوچک‌ترین و بزرگ‌ترین مقدار ممکن برای $\P(A\cap B)$ را بیابید و برای هر کدام یک مثال بسازید.
\answer{$0.7\le\P(A\cap B)\le0.8$}
\end{exercise}

\begin{exercise}
فرض کنید $A_1,A_2,\dots$ پیشامدهایی باشند که برای هر $n$، $\P(A_n)=1$. ثابت کنید $\P\bigl(\bigcap_{n=1}^\infty A_n\bigr)=1$. (راهنمایی: کران اجتماع را برای متمم‌ها به کار ببرید.) آیا نتیجه برای خانوادهٔ \emph{ناشمارا} هم درست است؟ (به نقطهٔ تصادفی در $[0,1]$ و پیشامدهای $A_x=[0,1]-\{x\}$ فکر کنید.)
\answer{برای خانوادهٔ ناشمارا نه: اشتراک همهٔ $A_x$ها تهی است.}
\end{exercise}

\begin{exercise}\hard
ثابت کنید نامساوی‌های بونفرونی زیر برقرارند:
\[ \sum_i\P(A_i)-\sum_{i<j}\P(A_i\cap A_j)\ \le\ \P\Bigl(\bigcup_{i=1}^nA_i\Bigr)\ \le\ \sum_i\P(A_i). \]
(راهنمایی: برای نامساوی سمت چپ از استقرا روی $n$ استفاده کنید.)
\end{exercise}

\begin{exercise}
یک عدد صحیح به تصادف از $\{1,2,\dots,1000\}$ انتخاب می‌شود. احتمال اینکه بر ۳ یا ۵ یا ۷ بخش‌پذیر باشد چقدر است؟
\answer{$543/1000$}
\end{exercise}

\begin{exercise}
یک تاس سالم ۶ بار پرتاب می‌شود. احتمال اینکه همهٔ اعداد ۱ تا ۶ ظاهر شوند چقدر است؟ احتمال اینکه دقیقاً دو عدد متمایز ظاهر شوند چطور؟
\answer{$6!/6^6\approx0.0154$ \quad؛\quad $\binom62(2^6-2)/6^6\approx0.0199$}
\end{exercise}

\begin{exercise}
در یک کلاس ۶۰ نفره، احتمال اینکه دست‌کم یک نفر روز تولدش با روز تولد \emph{شما} یکی باشد چقدر است؟ آن را با احتمال مسئلهٔ روز تولد برای $n=60$ مقایسه کنید.
\answer{$1-(364/365)^{60}\approx0.152$ در برابر $0.994$}
\end{exercise}

\begin{exercise}
یک کمیتهٔ ۴ نفره به تصادف از ۶ زوج (۱۲ نفر) انتخاب می‌شود. احتمال اینکه هیچ زوجی هر دو در کمیته نباشند چقدر است؟
\answer{$\binom64 2^4/\binom{12}4=240/495=16/33$}
\end{exercise}

\begin{exercise}[مثلث در گراف تصادفی]
\label{exr:triangle}
روی ۴ رأس $\{A,B,C,D\}$ هر یک از ۶ یالِ ممکن را با احتمال $p$ رسم می‌کنیم؛ دقیق‌تر، $\Omega$ مجموعهٔ همهٔ $2^6$ گراف ممکن است و گرافی با $m$ یال وزن $p^m(1-p)^{6-m}$ دارد.
\begin{enumerate}[label=(\alph*)]
  \item با قضیهٔ دوجمله‌ای نشان دهید اگر $S$ مجموعه‌ای مشخص از $s$ یال باشد، $\P(\fa{همهٔ یال‌های $S$ موجودند})=p^s$. (با $s=0$ نتیجه بگیرید که وزن‌ها جمعشان ۱ است.)
  \item چهار مثلث ممکن را $\Delta_1,\dots,\Delta_4$ بنامید. با اصل شمول و عدم شمول احتمال اینکه گراف \emph{دست‌کم یک مثلث} داشته باشد را بیابید. (راهنمایی: هر دو مثلث دقیقاً یک یال مشترک دارند و هر سه مثلث با هم همهٔ ۶ یال را می‌پوشانند.)
\end{enumerate}
\answer{(ب) $4p^3-6p^5+3p^6$ \ (بررسی: برای $p=1$ برابر ۱ است.)}
\end{exercise}

\begin{exercise}[بزرگ‌ترین خوشه در گراف تصادفی]\hard
در گراف تصادفی اردوش--رنی $G(n,p)$ هر دو رأس از $n$ رأس با احتمال $p$ به هم وصل‌اند (مدل دقیق مثل تمرین \ref{exr:triangle}: گرافی با $m$ یال وزن $p^m(1-p)^{\binom n2-m}$ دارد). یک \emph{$k$-خوشه} مجموعه‌ای از $k$ رأس است که همهٔ رأس‌هایش دوبه‌دو به هم وصل‌اند.
\begin{enumerate}[label=(\alph*)]
  \item با استفاده از تمرین \ref{exr:triangle}(الف) نشان دهید برای هر مجموعهٔ $k$رأسیِ مشخص $K$: \ $\P(\fa{$K$ خوشه است})=p^{k(k-1)/2}$.
  \item با کران اجتماع نشان دهید $\P(\fa{وجود یک $k$-خوشه})\le\binom nk p^{k(k-1)/2}\le n^k p^{k(k-1)/2}$. چرا محاسبهٔ دقیق این احتمال با شمول و عدم شمول عملاً ممکن نیست؟
  \item نشان دهید اگر $k\ge1+\dfrac{2(\ln n+\ln\frac1\delta)}{\ln\frac1p}$، این احتمال حداکثر $\delta$ است.
  \item برای $n=1000$، $p=\tfrac12$ و $\delta=0.01$ کوچک‌ترین چنین $k$ای را بیابید و نتیجه را تعبیر کنید.
\end{enumerate}
(این ایده، که اردوش در ۱۹۴۷ به کار برد، «روش احتمالاتی» نام دارد.)
\answer{(د) $k=35$: با احتمال دست‌کم ۹۹٪ هیچ گروه ۳۵ رأسیِ کاملاً به‌هم‌وصل وجود ندارد.}
\end{exercise}

